\chapter{Historical Hypotheses and Reconstruction Notes}\label{app:history}

\section{Historical note}\label{sec:historical-note}

The modern consonantal grades combine changes of primary articulation,
secondary palatalisation, and affrication. A possible history distinguishes
their sources rather than deriving all three from one sound change.
The developments below are hypotheses, with relative ordering specified
only where one development supplies the input to another. Modern contrasts
do not alone establish an ancestral hard--soft inventory; Wandl and
Kavitskaya's reconstruction of Common Slavic palatalisation illustrates the
need to reconstruct individual contrasts separately \cite{wandlkavitskaya2023}.

\noindent\textit{Glide absorption and lexical softness.}\quad
Earlier consonant--glide sequences could have yielded secondary palatality
on labials and changes of primary articulation on coronals, with different
segments reaching their modern outputs at different stages. Absorption of
the glide would precede the establishment of softness independent of an
adjacent /j/. Such a development could supply lexical palatality before
back vowels, as in modern \ird{pjašt-}. The digraph \orth{pj}
records its modern palatalised consonant without establishing that history.
The restricted lexical /kʲ/ of \ird{kěda} does not establish a former
/k/--/kʲ/ alternation or a historical source for its secondary palatality.
Its contrast with /c/ in \ird{těda} is synchronic; the two articulations
imply no particular ordering of earlier changes.

\noindent\textit{Velar grades and morphological retention.}\quad
The front-grade /k/--/ts/ and mixed-grade /k/--/tʃ/ relations could descend
from changes in different front-vocalic and consonant--glide environments.
Neither the order of these changes nor their intermediate articulations
is established by the modern alternants. Once changes to vowels or glides
obscured their conditioning, the resulting stem relations could have been
retained by particular suffix classes and extended analogically within
those classes. This would permit a later or differently inherited
front-vocalic exponent to remain plain-grade, as conditional \ird{-ic}
does, while other exponents retained mutation. The suffix-specific Czech
patterns provide a synchronic comparison for such differentiation
\cite{caviranivandenwyngaerd2026}, without establishing this Iridian chronology.

\noindent\textit{Convergent sibilant outputs.}\quad
The shared postalveolar outputs of coronal PAL, velar mixed selection, and antipassive
replacement are compatible with convergence of earlier articulations, but
the modern correspondences establish neither an earlier third sibilant
series nor the chronology of a merger.

\noindent\textit{Lateral vocalisation.}\quad
Earlier coda /l/ could have developed a dorsal articulation before losing
its lateral constriction and yielding [w], while a palatal counterpart
retained lateral articulation as [ʎ]. An intermediate velarised lateral is
a possible phonetic stage, not an independently reconstructed Iridian
segment. The modern [w]--[ʎ] distribution is compatible with this pathway,
but does not fix the historical order of vocalisation and the creation of
palatal laterals. Synchronically, stem selection and mutation precede
coda vocalisation (§\,\ref{sec:l-vocalisation}).

\noindent\textit{Construction-specific replacement.}\quad
Earlier fusion with an antipassive formative, followed by morphological
reanalysis and analogy, could have produced the replacement and additive
classes now associated with \ird{-aš-}. Their modern distribution does
not establish the original form of the formative or a single historical
palatalisation rule. Learned \ird{ř} likewise supplies no evidence for
a former productive rhotic mutation. These histories must remain separate
from the productive mappings until lexical correspondences establish
their connections.

\subsection{Vowel quality and the diphthong}

The high and mid short--long pairs differ in quality as well as quantity, although
their earlier phonetic values are not independently established. One possible
history would derive /ɪ, iː/, /ɛ, eː/, /ɔ, oː/, and /ʊ, uː/ from earlier
*/i, iː/, */e, eː/, */o, oː/, and */u, uː/, respectively, through lowering
of the short members, while */a, aː/ retained its quality. An earlier back
unrounded */ɯ, ɯː/ series could have centralised to /ɨ, ɨː/, with subsequent
lowering of short /ɨ/ yielding its modern [ɨ̞] target. These correspondences
constitute a hypothesis about the sources of the inventory; they do not
establish that the changes occurred together or determine their ordering
relative to the stem alternations below. Lexical correspondences, the
distribution of stable and alternating vowels, and the chronology of loans
would have to distinguish this history from other possible sources of the
same modern qualities.

The origin of /ou̯/ is likewise unresolved. An earlier vowel sequence and
diphthongisation of an earlier long vowel are possible sources, but neither
is established by the modern spelling \ird{ou}. Any derivation from an earlier
long /oː/ would also have to explain the distribution of the surviving
monophthong /oː/. Modern coda vocalisation supplies a separate, synchronically
productive source of vowel--[w] sequences and does not by itself account for
the phonemic diphthong.

\subsection{Old Iridian stem vocalism}\label{sec:old-iridian-vocalism}

The following tentative reconstruction explores how ordinary Old Iridian
short and long vowels, initial stress, and suffixal inflection might have
yielded the vocalic stem classes. Starred forms are hypothetical stages,
not independently attested forms; the proposed chronology remains to be
tested against the full lexical distribution. The reconstruction requires no front
and back yers and no phonologically empty vowels in modern representations.

\noindent\textit{Post-tonic syncope.}\quad
Short unstressed \ird{e} in the final stem syllable was lost before eligible
vowel-initial endings where consonant preservation allowed a licit
syllabification: hypothetical \ird{*Janek-u} $>$ \ird{*Jank-u} and
\ird{*perek-a} $>$ \ird{*perk-a}. These syncope steps alone do not derive
the strength of modern \ird{pérk-}. The corresponding unsuffixed forms retained the vowel in their
closed final syllable. Long unstressed vowels resisted this syncope; their
later shortening, together with later loans and morphological renewal, supplied
stable vowels in environments now shared with alternating stems. The modern unstressed long vowels have later lexical or morphological
sources, while the earlier quantity contrast survives indirectly in the
distinction between stable and alternating stems.

\noindent\textit{Metaphony.}\quad
After syncope, a high vowel in the next vocalic position could raise a stressed
mid vowel; consonants were transparent, while an intervening full vowel blocked
the relation. A hypothetical sequence \ird{*lobek-u} $>$ \ird{*lobk-u}
$>$ \ird{*lubk-u} could supply a raised stem, but not by itself the long
vowel of modern G3 \ird{lúbk-}. Height raising and the source of
morphological strength require separate explanations. High-suffix-triggered
raising of stressed vowels has a typological parallel in the Veneto pattern
analysed by Walker \cite{walker2001}.

\noindent\textit{Stem levelling and ending replacement.}\quad
On this hypothesis, high-vowel obliques could have supplied a raised bound
stem subsequently extended to non-high endings and anchoring. Unraised
reduced stems could instead have spread from older non-high endings, with
case-ending replacement obscuring the earlier conditioning. This offers a
possible account of class membership, not an established chronology for
\ird{pérk-}, \ird{lúbk-}, or \ird{Mercí}. In particular, modern G1 and G3
must be distinguished from height raising alone; a modern suffix vowel does
not determine either class membership or root quantity.

\noindent\textit{Verbal constructions.}\quad
For the irrealis, a sequence of ordinary \ird{*-i} and a following formative
\ird{*=ew} supplies a possible source: raising preceded glide formation and
jotation, as in hypothetical \ird{*zděk-i=ew} $>$ \ird{*zďik-i=ew}
$>$ \ird{*zďik-jew}. Glide absorption in \ird{k+j > č} and development
of the labial glide into \ird{v} could account for the consonants, while
the source of strength in modern \ird{zďíčev} remains a separate question.
Negative irrealis \ird{-ir} belongs to a different modern construction;
its mid root vowel does not establish the historical chronology.

The perfective--linker construction was integrated after nominal syncope had
ceased. Local suffixal assimilation gave \ird{*-ek-i > -ik-i}; the older
word-level linker was subsequently replaced by \ird{e}, reconstructed as a
phrase-level connective, while the bound perfective grade survived. This
suffixal history differs from stressed-root metaphony. The surviving suffixal correspondence is selected by the modern linker
construction. The participant nominaliser
\ird{-ou} selects the same bound perfective grade.

\noindent\textit{Quantity.}\quad
A separate hypothetical development could phonologise greater duration before
final sonorants and voiced fricatives in a subset of stems: reconstructed
\ird{*straz} $>$ \ird{stráz}, beside short-root
\ird{*straz-y}, and \ird{*trav} $>$ \ird{tráv}, beside \ird{trava}.
A phonetic-duration source distinct from lost-vowel compensation is developed
for Western Slavic by Scheer \cite{scheer2017}, providing a comparative model
for the reconstructed Iridian development.
Such a pathway would require subsequent morphological distribution and analogy
to account for modern plural strength and the G1/G3 stem classes. The modern
system does not establish a single origin for its quantity alternations;
independent lexical long vowels can remain long before suffixes.
Modern final devoicing and coda vocalisation obscure the earlier consonantal
conditioning: \ird{stráz} ends in [s], while \ird{tráv} ends in [w].



\section{Historical interpretation}

A possible history of the system begins with a small inherited or internally
developed set of bound directional and relational formatives, whose lexical
range subsequently expanded in the setting of prolonged Central European
contact described in Chapter~\ref{ch:overview}. On this hypothesis, outward
direction in \ird{le-} extends through emergence to extraction and disclosure,
while spatial entry in \ird{ra-} extends to conceptual penetration.
A circuitous path supplies the basis for circumvention and an indirect approach
in \ird{vra-}; separation of trajectories supplies dispersal and restricted
distributive extensions in \ird{ša-}. Spatial accumulation develops through
assembly and integration to synthesis in \ird{na-}, while return along an
earlier path supplies the model for restitution in \ird{ru-}. These are
hypothesised semantic pathways, not established etymologies or reconstructions
of ancestral forms.

The formal identity of accumulative \ird{na-} and locative postpositional
\ird{na} makes common origin a possibility, although their present meanings
and distributions do not establish the direction or intermediate stages of such
a development. Nor does the existence of a modern postposition provide a
necessary source for every preverb. Lexicalisation within Iridian and
contact-induced extension of an existing pattern remain separate historical
possibilities: a parallel with a neighbouring Slavic pattern would not by itself
establish borrowing of a formative or adoption of its aspectual grammar.
